\section{Man Makes himself Happy and Immortal}

\paragraph{Theory of the Will}

In \emph{The Key of the Mysteries}\footnote{This book was translated by Aleister Crowley, who regarded Levi as his spiritual predecessor in some sense.}, \textbf{Eliphas Levi}, who was a devout Catholic at least in his own mind, promises to harmonize reason with faith, and to

\begin{quotex}
exhibit true religion with such characters, that no one, believer or unbeliever, can fail to recognize it; that will be the absolute religion.
\end{quotex}

Part of this involves training of the will because

\begin{quotex}
Human life and its innumerable difficulties have for object, in the ordination of eternal wisdom, the education of the will of man.
\end{quotex}

These are the main points of the Theory of the Will, followed by our commentary.

\begin{enumerate}
\item The dignity of man consists in doing what he wills and in willing the good, in conformity with the knowledge of truth.
\item The good in conformity with the true, is the just.
\item Justice is the practice of reason.
\item Reason is the word of reality.
\item Reality is the science of truth.
\item Truth is idea identical with being.
\item Man arrives at the absolute idea of being by two roads, experience and hypothesis.
\item Experience is science and hypothesis is faith.
\end{enumerate}

\paragraph{Commentary}


\begin{enumerate}
\item Love means to will the good. So point (1) may be rephrased as ``Do what you will and Love, in conformity with the knowledge of Truth.'' So the dignity of man consists in loving and doing his true will. Corollary: Not every man has dignity.
\item In order to will the good, one must know what is the good, and that is the just. Justice is having and doing what is one's own.
\item According to Levy, Justice is the criterion of the truth of the Absolute religion in ethics, hence the exercise of the practical reason.
\item Reason is the criterion of the truth of the Absolute religion in judgment. Reason is the word, or Logos, of reality. Reason is the ability to make judgments, the aim of which is the finding of justice.
\item Reality is the criterion of the truth of the Absolute religion in science. Science knows reality by experience, whose object is the phenomenal world. Reason needs to be in conformity with reality. ``One has to know the extent of one's knowledge and ignorance when one makes a judgment.'' (Tomberg)
\item A fundamental metaphysical principle according to Guenon is that knowledge means the idea is identical with being, or, knowing equals being.
\item The absolute idea of being is Truth.
\item Science is Reason based on experience. The object of faith is Mystery.
\end{enumerate}

Levy provides us with two definitions.

\begin{description}

\item[God]

The supreme personification of law
\item[Liberty]

The right to do one's duty
\end{description}


This quote, then, acts as summary.

\begin{quotex}
As there is no liberty for man but in the order which results from the true and the good, one may say that the conquest of liberty is the great work of the human soul. \emph{Man, by freeing himself from his evil passions and their slavery, creates himself, as it were, a second time.} Nature made him living and suffering; he makes himself happy and immortal; he thus becomes the representative of divinity upon earth. And (relatively) exercises its almighty power.
\end{quotex}

\flrightit{Posted on 2010-10-03 by Cologero}

\begin{center}* * *\end{center}

\begin{footnotesize}
\begin{sffamily}
\texttt{Margi on 2010-10-04 at 10:51 said:}

Eliphus Levy's real name was Abbe Louis Constant. He said that the Bible and the Qu'ran are translations of the same book.

\hfill

\texttt{Heraklion on 2011-07-10 at 01:32 said:}

Eliphas Levi was definitely philosophically mixed-up, but he is not even remotely comparable to the cosmic Chthulian slime-essence of Crowley's cleverly-concealed, clinically-psychopathic personality. Eliphas Levi was good-intentioned but foolish; Crowley merely was hate and chaos personified vacuously... 

\hfill

\end{sffamily}
\end{footnotesize}

